% Options for packages loaded elsewhere
\PassOptionsToPackage{unicode}{hyperref}
\PassOptionsToPackage{hyphens}{url}
\documentclass[
  british,
  a4paper,
]{article}
\usepackage{xcolor}
\usepackage[margin=1in]{geometry}
\usepackage{amsmath,amssymb}
\setcounter{secnumdepth}{-\maxdimen} % remove section numbering
\usepackage{iftex}
\ifPDFTeX
  \usepackage[T1]{fontenc}
  \usepackage[utf8]{inputenc}
  \usepackage{textcomp} % provide euro and other symbols
\else % if luatex or xetex
  \usepackage{unicode-math} % this also loads fontspec
  \defaultfontfeatures{Scale=MatchLowercase}
  \defaultfontfeatures[\rmfamily]{Ligatures=TeX,Scale=1}
\fi
\usepackage{lmodern}
\ifPDFTeX\else
  % xetex/luatex font selection
  \setmainfont[]{Noto Serif}
  \setmonofont[]{Noto Sans Mono}
\fi
% Use upquote if available, for straight quotes in verbatim environments
\IfFileExists{upquote.sty}{\usepackage{upquote}}{}
\IfFileExists{microtype.sty}{% use microtype if available
  \usepackage[]{microtype}
  \UseMicrotypeSet[protrusion]{basicmath} % disable protrusion for tt fonts
}{}
\makeatletter
\@ifundefined{KOMAClassName}{% if non-KOMA class
  \IfFileExists{parskip.sty}{%
    \usepackage{parskip}
  }{% else
    \setlength{\parindent}{0pt}
    \setlength{\parskip}{6pt plus 2pt minus 1pt}}
}{% if KOMA class
  \KOMAoptions{parskip=half}}
\makeatother
\usepackage{color}
\usepackage{fancyvrb}
\newcommand{\VerbBar}{|}
\newcommand{\VERB}{\Verb[commandchars=\\\{\}]}
\DefineVerbatimEnvironment{Highlighting}{Verbatim}{commandchars=\\\{\}}
% Add ',fontsize=\small' for more characters per line
\newenvironment{Shaded}{}{}
\newcommand{\AlertTok}[1]{\textcolor[rgb]{1.00,0.00,0.00}{\textbf{#1}}}
\newcommand{\AnnotationTok}[1]{\textcolor[rgb]{0.38,0.63,0.69}{\textbf{\textit{#1}}}}
\newcommand{\AttributeTok}[1]{\textcolor[rgb]{0.49,0.56,0.16}{#1}}
\newcommand{\BaseNTok}[1]{\textcolor[rgb]{0.25,0.63,0.44}{#1}}
\newcommand{\BuiltInTok}[1]{\textcolor[rgb]{0.00,0.50,0.00}{#1}}
\newcommand{\CharTok}[1]{\textcolor[rgb]{0.25,0.44,0.63}{#1}}
\newcommand{\CommentTok}[1]{\textcolor[rgb]{0.38,0.63,0.69}{\textit{#1}}}
\newcommand{\CommentVarTok}[1]{\textcolor[rgb]{0.38,0.63,0.69}{\textbf{\textit{#1}}}}
\newcommand{\ConstantTok}[1]{\textcolor[rgb]{0.53,0.00,0.00}{#1}}
\newcommand{\ControlFlowTok}[1]{\textcolor[rgb]{0.00,0.44,0.13}{\textbf{#1}}}
\newcommand{\DataTypeTok}[1]{\textcolor[rgb]{0.56,0.13,0.00}{#1}}
\newcommand{\DecValTok}[1]{\textcolor[rgb]{0.25,0.63,0.44}{#1}}
\newcommand{\DocumentationTok}[1]{\textcolor[rgb]{0.73,0.13,0.13}{\textit{#1}}}
\newcommand{\ErrorTok}[1]{\textcolor[rgb]{1.00,0.00,0.00}{\textbf{#1}}}
\newcommand{\ExtensionTok}[1]{#1}
\newcommand{\FloatTok}[1]{\textcolor[rgb]{0.25,0.63,0.44}{#1}}
\newcommand{\FunctionTok}[1]{\textcolor[rgb]{0.02,0.16,0.49}{#1}}
\newcommand{\ImportTok}[1]{\textcolor[rgb]{0.00,0.50,0.00}{\textbf{#1}}}
\newcommand{\InformationTok}[1]{\textcolor[rgb]{0.38,0.63,0.69}{\textbf{\textit{#1}}}}
\newcommand{\KeywordTok}[1]{\textcolor[rgb]{0.00,0.44,0.13}{\textbf{#1}}}
\newcommand{\NormalTok}[1]{#1}
\newcommand{\OperatorTok}[1]{\textcolor[rgb]{0.40,0.40,0.40}{#1}}
\newcommand{\OtherTok}[1]{\textcolor[rgb]{0.00,0.44,0.13}{#1}}
\newcommand{\PreprocessorTok}[1]{\textcolor[rgb]{0.74,0.48,0.00}{#1}}
\newcommand{\RegionMarkerTok}[1]{#1}
\newcommand{\SpecialCharTok}[1]{\textcolor[rgb]{0.25,0.44,0.63}{#1}}
\newcommand{\SpecialStringTok}[1]{\textcolor[rgb]{0.73,0.40,0.53}{#1}}
\newcommand{\StringTok}[1]{\textcolor[rgb]{0.25,0.44,0.63}{#1}}
\newcommand{\VariableTok}[1]{\textcolor[rgb]{0.10,0.09,0.49}{#1}}
\newcommand{\VerbatimStringTok}[1]{\textcolor[rgb]{0.25,0.44,0.63}{#1}}
\newcommand{\WarningTok}[1]{\textcolor[rgb]{0.38,0.63,0.69}{\textbf{\textit{#1}}}}
% definitions for citeproc citations
\NewDocumentCommand\citeproctext{}{}
\NewDocumentCommand\citeproc{mm}{%
  \begingroup\def\citeproctext{#2}\cite{#1}\endgroup}
\makeatletter
 % allow citations to break across lines
 \let\@cite@ofmt\@firstofone
 % avoid brackets around text for \cite:
 \def\@biblabel#1{}
 \def\@cite#1#2{{#1\if@tempswa , #2\fi}}
\makeatother
\newlength{\cslhangindent}
\setlength{\cslhangindent}{1.5em}
\newlength{\csllabelwidth}
\setlength{\csllabelwidth}{3em}
\newenvironment{CSLReferences}[2] % #1 hanging-indent, #2 entry-spacing
 {\begin{list}{}{%
  \setlength{\itemindent}{0pt}
  \setlength{\leftmargin}{0pt}
  \setlength{\parsep}{0pt}
  % turn on hanging indent if param 1 is 1
  \ifodd #1
   \setlength{\leftmargin}{\cslhangindent}
   \setlength{\itemindent}{-1\cslhangindent}
  \fi
  % set entry spacing
  \setlength{\itemsep}{#2\baselineskip}}}
 {\end{list}}
\usepackage{calc}
\newcommand{\CSLBlock}[1]{\hfill\break\parbox[t]{\linewidth}{\strut\ignorespaces#1\strut}}
\newcommand{\CSLLeftMargin}[1]{\parbox[t]{\csllabelwidth}{\strut#1\strut}}
\newcommand{\CSLRightInline}[1]{\parbox[t]{\linewidth - \csllabelwidth}{\strut#1\strut}}
\newcommand{\CSLIndent}[1]{\hspace{\cslhangindent}#1}
\ifLuaTeX
\usepackage[bidi=basic,shorthands=off]{babel}
\else
\usepackage[bidi=default,shorthands=off]{babel}
\fi
\ifPDFTeX
\else
\babelfont{rm}[]{Noto Serif}
\fi
\ifLuaTeX
  \usepackage{selnolig} % disable illegal ligatures
\fi
\setlength{\emergencystretch}{3em} % prevent overfull lines
\providecommand{\tightlist}{%
  \setlength{\itemsep}{0pt}\setlength{\parskip}{0pt}}
\usepackage{array}
\usepackage{tabularx}
\usepackage{bookmark}
\IfFileExists{xurl.sty}{\usepackage{xurl}}{} % add URL line breaks if available
\urlstyle{same}
\hypersetup{
  pdflang={en-GB},
  hidelinks,
  pdfcreator={LaTeX via pandoc}}

\author{}
\date{}

\begin{document}

\section{Sound-change reports, alpha
01}\label{sound-change-reports-alpha-01}

\emph{Generated from
\texttt{Germanic/docs/sound\_changes/change\_reports/report\_manifest.tsv}.
This assembly includes only manifest-backed \texttt{pilot} and
\texttt{full} production reports.}

\subsection{Early vocalic/final
corridor}\label{early-vocalicfinal-corridor}

\subsubsection{Sound-change report}\label{sound-change-report}

\paragraph{Historical formulation}\label{historical-formulation}

This entry treats \textbf{SC016, SC017, SC019, and SC020} as one
corridor because the four changes interact in actual derivations, not
because the handbooks present them as a ready-made chapter. The corridor
begins with the West Saxon \texttt{geoc\ /\ geong\ /\ geoguþ}
phenomenon, passes through the standard lowering of Germanic \texttt{u}
before a following non-high vowel, turns into final-vowel structure
through unstressed \texttt{*-ō\ \textgreater{}\ *-u}, and closes with
West Germanic loss of final \texttt{*z} before rhotacism
(\citeproc{ref-Kaluza1906}{Kaluza 1906, sect. 90};
\citeproc{ref-Bulbring1902}{Bülbring 1902, sects 298--299};
\citeproc{ref-Campbell1959}{Campbell 1959, sect. 44};
\citeproc{ref-Campbell1959}{1959, sect. 115};
\citeproc{ref-RingeTaylor2014}{Ringe and Taylor 2014, 267};
\citeproc{ref-Crist2002}{Crist 2002}).

The sequence is therefore historical but uneven. SC017 is the clearest
standard handbook sound change. SC016 is real, but the literature
usually presents it as a West Saxon palatal-glide or rising-diphthong
phenomenon rather than as a widely named law
(\citeproc{ref-Campbell1959}{Campbell 1959, sect. 44}). SC019 is
historical, but the literature usually frames it through \texttt{ō}-stem
and final-vowel behavior rather than under the CAPR label \emph{NWGmc
Final Long O Raising} (\citeproc{ref-RingeTaylor2014}{Ringe and Taylor
2014, 267}; \citeproc{ref-Fulk2018}{Fulk 2018}). SC020 is likewise
historical, but the clearest formulation is the loss of word-final
\texttt{*z} in unstressed syllables with major morphological
consequences, not the exact one-line FOMA rule used by CAPR
(\citeproc{ref-Crist2001}{Crist 2001}, \citeproc{ref-Crist2002}{2002};
\citeproc{ref-Kilday2024}{Kilday 2024}).

\paragraph{Source tradition}\label{source-tradition}

The source tradition begins at the left edge with the older Old English
grammarians. Kaluza and Bülbring treat \texttt{geoc}, \texttt{geong},
and \texttt{geoguþ} as part of the West Saxon reflexes of initial
Germanic \texttt{j}, while Campbell gives the most usable short
formulation for book prose: the spellings in forms such as \texttt{geoc}
reflect rising diphthongs formed when palatal glides developed before
back vowels (\citeproc{ref-Kaluza1906}{Kaluza 1906, sect. 90};
\citeproc{ref-Bulbring1902}{Bülbring 1902, sects 298--299};
\citeproc{ref-Campbell1959}{Campbell 1959, sect. 44}). That is enough to
justify SC016 as a real phenomenon, but it also shows why the subsection
should be phrased cautiously: the sources describe a real West Saxon
development more often than they name a discrete sound law.

For SC017 the handbook tradition is firmer. Kaluza states the rule early
and compactly, Campbell gives the standard textbook wording
``\texttt{u\ \textgreater{}\ o\ before\ mid\ and\ low\ vowels},'' and
Fulk provides the clearest modern formulation with the relevant blocking
conditions (\citeproc{ref-Kaluza1906}{Kaluza 1906, sect. 66};
\citeproc{ref-Campbell1959}{Campbell 1959, sect. 115};
\citeproc{ref-Fulk2018}{Fulk 2018, sect. 4.3}). Luick is especially
useful here as a cautionary witness, because he shows that the outputs
could later fluctuate analogically, as in the \texttt{scufel\ /\ scofel}
type (\citeproc{ref-Luick1914}{Luick 1914-\/-1940}). The corridor should
therefore treat SC017 as the strongest standard change in the set, while
still recognizing that the distribution of actual forms is not
mechanically uniform.

The tradition for SC019 is more ending-centered. Luick treats final
\texttt{ō\ \textgreater{}\ u} through final-vowel shortening, Ringe and
Taylor give the clearest concise statement that Proto-Germanic
word-final non-nasalized long \texttt{*-ō} became unstressed
\texttt{*-u} in Proto-Northwest-Germanic, and Fulk translates that
comparative claim into the \texttt{ō}-stem and light-stem endings
familiar from Old English (\citeproc{ref-Luick1914}{Luick 1914-\/-1940};
\citeproc{ref-RingeTaylor2014}{Ringe and Taylor 2014, 267};
\citeproc{ref-Fulk2018}{Fulk 2018}). That makes SC019 a historical
change with strong comparative support, but not one whose CAPR label
should necessarily become the prose heading unchanged.

SC020 continues the same move from root vocalism into endings and weak
tails. Campbell and Hogg preserve the textbook frame that Germanic
\texttt{z} is later lost or rhotacized, while Crist's dissertation and
2002 handout formulate the key historical claim more sharply as West
Germanic loss of word-final \texttt{*z} in unstressed syllables before
rhotacism; Kilday's recent summary treats that core development as
already established background (\citeproc{ref-Campbell1959}{Campbell
1959}; \citeproc{ref-Hogg1992}{Hogg 1992};
\citeproc{ref-Crist2001}{Crist 2001}, \citeproc{ref-Crist2002}{2002};
\citeproc{ref-Kilday2024}{Kilday 2024}). This means the right edge of
the corridor is also historically real, but its literature is more
morphology-heavy than the bare CAPR rule suggests.

\paragraph{CAPR implementation}\label{capr-implementation}

CAPR turns those historical phenomena into four explicit ordered stages.
SC016 formalizes the West Saxon palatal-glide / rising-diphthong
behavior by rewriting initial palatal + \texttt{u/ú} sequences to
palatal + \texttt{eu/éu}, extended across the initial palatal inventory
used in the model. SC017 is the most ordinary sound-law implementation
of the four: it lowers stressed \texttt{u/ú} to \texttt{o/ó} before a
following non-high vowel while explicitly blocking the relevant nasal
and \texttt{j} environments. SC019 formalizes final unstressed
\texttt{*-ō\ \textgreater{}\ *-u} as a separate stage, while the
adjacent monosyllabic \texttt{ō} rule is split off into SC018. SC020
then deletes final \texttt{*z} outright with the simple rewrite
\texttt{{[}\{*z\}\ -\textgreater{}\ 0\ \textbar{}\textbar{}\ \_\ .\#.{]}},
even though the literature states the historical core more narrowly as
loss of word-final \texttt{*z} in unstressed syllables
(\citeproc{ref-Campbell1959}{Campbell 1959};
\citeproc{ref-Fulk2018}{Fulk 2018}; \citeproc{ref-RingeTaylor2014}{Ringe
and Taylor 2014}; \citeproc{ref-Crist2002}{Crist 2002}).

That explicitness is useful, but it is not neutral. SC016 is more
sharply bounded in CAPR than in the grammars. SC017 is closest to a
direct handbook implementation. SC019 isolates a literature-backed
ending development as a single stage with exact structural guards. SC020
is the strongest case where the formal model is broader than the
cleanest handbook formulation.

\paragraph{Place in the cascade}\label{place-in-the-cascade}

The relevant cascade slice is:

\begin{Shaded}
\begin{Highlighting}[]
\NormalTok{.o. NWGmcUnstressedAiMonophthongization}
\NormalTok{.o. NWGmcILowering}
\NormalTok{.o. OEWsPalatalGlide}
\NormalTok{.o. NWGmcULowering}
\NormalTok{.o. NWGmcStressedMonosyllableORaising}
\NormalTok{.o. NWGmcFinalLongORaising}
\NormalTok{.o. PGmcFinalZDeletion}
\end{Highlighting}
\end{Shaded}

This is the local neighborhood that makes the corridor visible. SC016,
SC017, SC019, and SC020 form the meaningful derivational spine. SC018
remains adjacent in the implementation, but its chronology card is
negative / boundary-limited, so it is not chapter-forming in the same
way. SC020 also points forward: once final \texttt{*z} is gone, the
model is already moving into the later final-syllable and weak-tail
territory that becomes more prominent in subsequent chapters.

\paragraph{Order evidence}\label{order-evidence}

The chronology cards recover three local order-sensitive links. First,
\texttt{SC016\ \textless{}\ SC017} is visible in \textbf{yoke}: PGmc
\texttt{*júką} yields expected OE \texttt{ġeoc}, but the wrong ordering
produces \texttt{ġoc}. The literature independently supports both the
\texttt{geoc} phenomenon and the \texttt{u\ \textgreater{}\ o} rule, but
not this exact local ordering statement; that is the contribution of
CAPR's order testing.

Second, \texttt{SC017\ \textless{}\ SC019} is visible in \textbf{nose},
\textbf{shovel}, and \textbf{sorrow}. If \texttt{u}-lowering is delayed
too long, the lowered root vowel never appears and the derivations come
out as \texttt{nusu}, \texttt{sċufl}, and \texttt{surg} instead of
\texttt{nosu}, \texttt{sċofl}, and \texttt{sorg}. This boundary is
literature-compatible, especially because Fulk places
\texttt{u}-lowering before the final \texttt{*-ō\ \textgreater{}\ *-u}
development (\citeproc{ref-Fulk2018}{Fulk 2018, sect. 4.3}), but the
exact local three-word test remains a CAPR result. \texttt{Nose} should
also be treated cautiously in polished prose because Kroonen gives the
etymon as \texttt{*nasō-\ \textasciitilde{}\ *nusō-} rather than as one
uncontested preform (\citeproc{ref-Kroonen2013}{Kroonen 2013}).

Third, \texttt{SC019\ \textless{}\ SC020} is visible in \textbf{rest}.
If final \texttt{*z} is deleted too early, final \texttt{*-ō} never
receives the treatment it needs and PGmc \texttt{*rástōz} comes out as
\texttt{rast} instead of expected \texttt{ræste}. The literature
strongly supports both ingredients of this interaction, but it does not
itself formulate the \texttt{rest} boundary as a local chronology
theorem. Here again CAPR supplies the explicit derivational hinge.

SC020 also has later broad/far links, especially toward SC040, but those
should be treated only as forward-looking context. They show that the
corridor leads into the later weak-tail system; they do not justify
expanding this entry into an undifferentiated graph chapter.

\paragraph{Interpretation}\label{interpretation}

The book-level value of this corridor is methodological as well as
historical. The literature gives four real phenomena, but it gives them
unevenly: one is a strong standard sound law, one is a real but weakly
named West Saxon sub-rule, and two belong most naturally to ending
structure and final-syllable history. The CAPR cascade then forces those
phenomena into explicit ordered rule form. The chronology cards,
finally, show where specific derivations break if that order is
disturbed.

That combination is what makes the corridor worth a production report.
It shows how the sound-change half of the book can move from scholarship
to model to controlled inference without confusing the three. The
handbooks do not give a ready-made four-part chapter, but they do
provide the historical material from which the corridor can be written
responsibly.

\paragraph{Remaining cautions}\label{remaining-cautions}

SC016 should probably be described in finished prose as a
\textbf{modeled West Saxon palatal-glide / rising-diphthong sub-rule}
rather than as a universally named sound law. SC019's label may likewise
need softening toward \textbf{final unstressed
\texttt{*-ō\ \textgreater{}\ *-u}} language. SC020's literature-backed
core is narrower than CAPR's unconditioned final-\texttt{z} deletion
rule, so the prose should keep the handbook formulation and the model
formulation distinct (\citeproc{ref-Crist2002}{Crist 2002};
\citeproc{ref-Kilday2024}{Kilday 2024}). \texttt{Nose} should continue
to carry an ablaut caution because of Kroonen's
\texttt{*nasō-\ \textasciitilde{}\ *nusō-} reconstruction
(\citeproc{ref-Kroonen2013}{Kroonen 2013}). Finally, the later SC020
links to SC041, SC042, and SC054 should remain background orientation
only; they should not expand this entry into a graph-driven chapter.

\protect\phantomsection\label{refs}
\begin{CSLReferences}{1}{1}
\bibitem[\citeproctext]{ref-Bulbring1902}
Bülbring, Karl D. 1902. \emph{Altenglisches Elementarbuch}. Winter.

\bibitem[\citeproctext]{ref-Campbell1959}
Campbell, Alistair. 1959. \emph{Old {English} Grammar}. Clarendon Press.

\bibitem[\citeproctext]{ref-Crist2001}
Crist, Sean J. 2001. {`Conspiracy in Historical Phonology'}. Ph.\{D\}.
dissertation, University of Pennsylvania.

\bibitem[\citeproctext]{ref-Crist2002}
Crist, Sean J. 2002. {`An Analysis of *{z} Loss in {West Germanic}'}.
Unpublished manuscript.

\bibitem[\citeproctext]{ref-Fulk2018}
Fulk, R. D. 2018. \emph{A Comparative Grammar of the Early {Germanic}
Languages}. Vol. 3. Studies in Germanic Linguistics. John Benjamins.

\bibitem[\citeproctext]{ref-Hogg1992}
Hogg, Richard M. 1992. \emph{A Grammar of Old English. Volume 1:
Phonology}. Blackwell.

\bibitem[\citeproctext]{ref-Kaluza1906}
Kaluza, Max. 1906. \emph{Historische Grammatik Der Englischen Sprache}.
2 vols. E. Felber.

\bibitem[\citeproctext]{ref-Kilday2024}
Kilday, Douglas G. 2024. {`{Crist's} Law, {Smith's} Law, and {English}
\emph{Wizen}'}. Unpublished manuscript.

\bibitem[\citeproctext]{ref-Kroonen2013}
Kroonen, Guus. 2013. \emph{Etymological Dictionary of {Proto-Germanic}}.
Vol. 11. Leiden Indo-European Etymological Dictionary Series. Brill.

\bibitem[\citeproctext]{ref-Luick1914}
Luick, Karl. 1914-\/-1940. \emph{Historische Grammatik Der Englischen
Sprache}. Tauchnitz.

\bibitem[\citeproctext]{ref-RingeTaylor2014}
Ringe, Don, and Ann Taylor. 2014. \emph{The Development of {Old
English}}. Vol. 2. A Linguistic History of English. Oxford University
Press.

\end{CSLReferences}

\end{document}
